\originalpreface
微观动力学保留可逆的粒子轨道，宏观气体却表现黏性、热传导与熵增；临界体系的相关长度又可远大于粒子尺度。高级统计力学的任务是说明这些描述怎样相容，并判断近似在哪个尺度有效。本册读者应已熟悉本科热力学、微正则/正则分布、理想量子气体与基本量子力学。仓库既有《热力学与统计物理》提供这些先修及8.333第15--22讲的部分补充；本卷用必要的符号连接它，重点展开动力学、关联与统计场论，不重复本科专题的全部例题。

课程主线为Mehran Kardar的\href{https://ocw.mit.edu/courses/8-333-statistical-mechanics-i-statistical-mechanics-of-particles-fall-2013/}{8.333，Fall 2013}及\href{https://ocw.mit.edu/courses/8-334-statistical-mechanics-ii-statistical-physics-of-fields-spring-2014/}{8.334，Spring 2014}。两站各公开26份课历对应讲义，官方都说明讲义来自更早版本；若文件未给确切原稿年份，本稿不猜造年份。课历的2013、2014日期、早期讲义版本与往年复习题年代分开记录。每门课有6套公开作业。8.333考试页列3套阶段复习与一套期末复习，8.334列3套阶段复习，都有公开官方解；8.333前三包标Fall2007，期末包含1998、1999、2000、2003--2006历史卷；8.334三包标Spring2014复习题，不能把复习题包称为该学期实际测试卷。8.333成绩为作业30\%、三测验45\%、期末25\%；8.334为作业30\%、三测验45\%、最终项目25\%。没有取得2013、2014当期实际考试卷及学生项目，均不补造。

第一编连接粒子、系综与流体，第二编系统建立场、相变、尺度与RG，第三编收录实际公开问题及中文解答。无公开答案的作业标为编者独立补写；复习包注明可获取官方解；中文解由编者独立推导，并在题内单列官方对照和确定可核验的修正。旧卷已写的8.333作业5仍核对原题并给固定前置定位，重复问题采用明确回指，而不作为新增解答计数。视频各26讲保留官方链接，未宣称逐分钟转录全部视频。

来源署名Mehran Kardar、MIT OpenCourseWare及文件所列作者。本项目作中文翻译、重编、推导补充与自绘矢量图，基于OCW部分采用\href{https://creativecommons.org/licenses/by-nc-sa/4.0/}{CC BY-NC-SA 4.0}，遵守署名、非商业及相同方式共享，并保留逐件例外；本项目不代表MIT认可。商业参考书只留出版信息与原题号，不打包整书或以答案反造书本题面。课程公开目录、实际文件与完成范围分别记录在学科来源清单与卷末。
